\begin{titlepage}
\centering
\vspace*{1.2cm}
{\Large\bfseries Undergraduate Thesis\par}
\vspace{0.6cm}
{\large School of Mathematics\par}
\vfill
{\LARGE\bfseries On the Sum of the Three Largest\\[0.35cm] Adjacency Eigenvalues of a Graph\par}
\vspace{1.0cm}
{\large A joint moment bound, the structure of extremal graphons,\\ and numerical verification\par}
\vfill
{\large Haoyang Ji\par}
\vspace{0.4cm}
{\normalsize 4 November 2026\par}
\vfill
\rule{0.55\textwidth}{0.4pt}\par
\vspace{0.3cm}
{\small This thesis expands the preprint \emph{A joint moment bound for the three
largest graph eigenvalues} (H.~Ji, 3 November 2026). All constants quoted
here were recomputed in exact arithmetic; the labels \emph{proved},
\emph{verified}, \emph{numerical} and \emph{gap} are used throughout to
separate the four kinds of statement.\par}
\end{titlepage}

\chapter*{Declaration}
\addcontentsline{toc}{chapter}{Declaration}
The work presented in this thesis is my own, carried out between August and
November 2026. Where a result, an idea, a construction or a constant is due
to another author, it is cited at the point of use. Two instances deserve to
be singled out here because they concern priority rather than technique.

\begin{itemize}[leftmargin=1.6em,itemsep=3pt]
\item The lower-bound constant $\cstar=(9+\sqrt5+2\sqrt{35})/18$ and the
optimising weight of the $K_1\vee C_5$ construction were announced earlier,
in conditional form, in the self-archived Zenodo deposits of
Newton~\cite{Newton1,Newton2}. No priority is claimed for that constant;
see Remark~\ref{rem-newton}.
\item The observation that the Ky Fan transfer argument needs only a zero
diagonal and off-diagonal entries in $[0,1]$ is already implicit in the
proof of~\cite[Proposition 6.1]{STW} and in the matrix setting
of~\cite{Mohar}. Proposition~\ref{weighted-form} records the resulting
weighted form of the main bound but claims no novelty for it.
\end{itemize}

Every displayed identity, factorisation and positivity assertion reproduced
from the preprint was re-checked in exact rational arithmetic and in the
quadratic fields $\mathbb Q(\sqrt3)$ and $\mathbb Q(\sqrt5)$, never in
floating point. Numerical experiments are reported as numerical evidence
only and are never used as a step in a proof.

\vspace{0.6cm}
\noindent Haoyang Ji \hfill 4 November 2026

\chapter*{Abstract}
\addcontentsline{toc}{chapter}{Abstract}
Let $\lambda_1(G)\ge\lambda_2(G)\ge\cdots$ be the adjacency eigenvalues of a
finite simple graph $G$ and write
$\Lambda_3(G)=\lambda_1(G)+\lambda_2(G)+\lambda_3(G)$. This thesis proves
\[
 \Lambda_3(G)\;\le\;c\,n-3\qquad(n\ge3),
\]
where
\[
 c=\frac{106106902139519}{107310156800000}+\frac{41\sqrt3}{240}
   -\frac{53}{9000000}=1.284673255723918505\ldots ,
\]
which improves the coefficient $(3+\sqrt5)/4=1.309016994\ldots$ obtained by
Sivashankar, Tang and Wakhare from the rank-three absolute projection
constant $\gamma(3)=(1+\sqrt5)/2$. The constant term $-3$ is unchanged.

The improvement comes from replacing two separate estimates by one joint
estimate. Writing $A=\E|X\cdot X'|$ and $B=|\E X|^2$ for an isotropic law
$\E XX^T=I_3$ on $\R^3$ with independent copy $X'$, the classical bounds
$A\le(1+\sqrt5)/2$ and $B\le1$ give exactly the coefficient
$(3+\sqrt5)/4$. We prove instead
$A+B\le 2c$, by an exact identity that couples $A$ and $B$ to the mean
deficit, the radial variance, the regression error and the weighted harmonic
moments of degrees three and four through a signed quartic majorant. All
error terms are retained in the identity, so the certificate is quantitative
rather than asymptotic.

Structurally, we prove that the canonical eigenfunction law of every
maximising graphon for the variational quantity
$\mathcal M=\sup_W\Sigma_3(W)$ has at most $24$ support points, that the
maximiser is almost everywhere a binary step kernel determined by the signs
of the dot products of those points, and that every positive clique in the
support consists of at most six linearly independent types. The argument
combines a second-variation inequality on the actual support, propagation of
moment relations, a dimension descent over sign regions, and a separating
family of binary observations.

The bound is shown to hold verbatim for $[0,1]$-weighted loopless graphs, and
weights are proved not to improve the asymptotic constant. Extensive
numerical verification is reported: all $2\,131\,016$ labelled graphs of orders
$3\le n\le7$ satisfy the bound, the unique extremal graph at $n=6$ is
$K_1\vee C_5$, and blow-ups of $K_1\vee C_5$ attain
$\cstar n-3+o(1)$ with $\cstar=1.281568196872168\ldots$.

Neither the optimal coefficient nor the extremal graphs are determined. The
residual gap $c-\cstar=0.003105058851750628\ldots$ and the structure of
maximisers are left open and are discussed in the final chapter.

\vspace{0.4cm}
\noindent\textbf{Keywords.} Adjacency eigenvalues; spectral sums; Ky Fan
maximum principle; graphons; absolute projection constants; isotropic
measures; extremal graph theory.

\chapter*{Acknowledgements}
\addcontentsline{toc}{chapter}{Acknowledgements}
I thank my supervisor for the patience and the many sharp questions that
shaped this work, and the authors of~\cite{STW,KLMPT,KMMP,Paulin,HW} for
making their preprints available so quickly, which allowed the present
argument to be positioned honestly against the current state of the
literature. I am also grateful to the authors of the recent preprints on spectral sums
for making their work publicly available so promptly, which allowed the
present argument to be positioned honestly against the current literature.

\clearpage
\thispagestyle{empty}
\null
\vfill
